\section{Background: The Networked XR Loop}\label{sec:background}

This section introduces the concepts that the rest of the survey relies on. Readers familiar with XR
systems may skip to \cref{sec:framework}.

\subsection{From standalone to networked XR}
XR spans a continuum. Virtual reality (VR) replaces the user's view with a synthetic one. Augmented
reality (AR) overlays virtual objects on the physical world, and mixed reality anchors them in it so
that they respond to physical surfaces. Each device tracks the user's head, and often the hands and eyes,
and redraws the scene for every new pose. Two levels of freedom matter for networking. With \emph{three
degrees of freedom} (3DoF) the user can only rotate the head, as in a 360-degree video. With \emph{six
degrees of freedom} (6DoF) the user can also move through the scene, as in a game or a volumetric
capture. A 3DoF scene can be prepared in advance as a sphere of pixels. A 6DoF scene must be rendered or
synthesised for each viewpoint, so it needs either geometry on the device or rendering close to it.

A standalone headset performs the whole loop locally: it senses a pose, renders, and displays the result.
Three pressures push parts of this loop into the network. First, \emph{content} may be too large to
store, or may be produced live elsewhere, as with 360-degree broadcasts and volumetric captures. Second,
\emph{computation} may exceed what a light, battery-powered device can sustain, as with photorealistic
rendering and AR scene understanding \cite{W3133913697}. Third, \emph{other participants}, whether
people or machines, are by definition remote \cite{W2561715836}.

\subsection{The motion-to-photon deadline}
The defining constraint of XR is the \emph{motion-to-photon} (MTP) latency: the time from a head movement
to the photons that show its consequence. When MTP exceeds roughly 15--20~ms the displayed world appears
to lag or swim, which degrades presence and can cause simulator sickness. That budget leaves roughly
14~ms for computation and communication together \cite{W2784711283}. Three properties of this deadline
shape every design in this survey. First, it is \emph{per frame}: at 72--120 frames per second, each
frame has its own deadline, and a late frame is a lost frame. Second, it applies to \emph{rotation} most
strictly. Local reprojection can correct a slightly stale frame for a new head orientation, so the
network-facing deadline for full interaction is looser than MTP, although it is still tight. Third, its
\emph{tail} matters more than its mean, because occasional late frames are what users notice.

\subsection{What XR traffic looks like}
The traffic of a networked XR session differs from video streaming or web traffic in four ways
(\cref{tab:background}). It is \emph{periodic and bursty}: rendered or decoded frames arrive at the
display rate, each a burst of packets with its own deadline. It is \emph{bidirectional}: pose, controller
input, camera frames and captured media flow upstream continuously, and the downstream content depends on
them. It is \emph{viewport-dependent}: the user sees only a fraction of a 360-degree or volumetric scene
at any moment, so the useful content depends on where the user looks. And it is often \emph{shared}:
several users' updates must be reconciled into one consistent world. The first two properties make XR a
real-time control problem rather than a throughput problem. The third creates the opportunity to send
less. The fourth makes the network responsible for consistency as well as delivery.

\begin{table}[t]
\centering\small
\caption{How networked XR traffic differs from conventional video streaming. Each row names a
property, its consequence for the network, and the section that treats it.}\label{tab:background}
\setlength{\tabcolsep}{4pt}
\begin{tabularx}{\linewidth}{@{}>{\raggedright\arraybackslash}p{2.3cm}>{\raggedright\arraybackslash}X>{\raggedright\arraybackslash}X>{\raggedright\arraybackslash}p{1.9cm}@{}}
\toprule
Property & Video streaming & Networked XR & Treated in\\
\midrule
Deadline & Seconds of buffer absorb delay & Each frame must arrive within a few tens of milliseconds & \cref{sec:compute,sec:provision}\\
Direction & Downlink, with small requests upstream & Continuous uplink of pose, input, camera and capture & \cref{sec:compute,sec:connect}\\
Viewport & The whole frame is shown & Only the viewport is seen, and it moves & \cref{sec:anticipate}\\
Representation & 2D video frames & Spherical video, point clouds, meshes, neural fields, haptic signals & \cref{sec:encode}\\
Participants & Independent viewers & Shared state, avatars, remote machines & \cref{sec:connect}\\
Quality & Rebuffering, bitrate, video quality & Latency, judder, sickness, presence, task success & \cref{sec:assure}\\
\bottomrule
\end{tabularx}
\end{table}

\subsection{Where the loop can run}
A networked XR system distributes its loop over four tiers: the \emph{device}, the \emph{access network}
(Wi-Fi or cellular), the \emph{edge} (servers at or near the access network) and the \emph{cloud}. Moving
a function to a farther tier gains resources and loses time. An edge server one wireless hop away adds a
few milliseconds of round-trip time, while a distant cloud data centre adds tens of milliseconds
\cite{W2805833065}. Mobile edge computing, as standardised for 5G, exists largely to put rendering,
caching and session state within the XR deadline \cite{W2902319104}. The access network is usually the
least predictable part of the path. Wireless capacity, delay and loss vary with distance, blockage,
interference and the number of users, and a directional millimetre-wave beam must follow a head that
turns quickly. The framework in the next section builds on these four tiers.
